\documentclass[aspectratio=169,11pt]{beamer}
\usetheme[
  progressbar=frametitle,
  sectionpage=none
]{moloch}
% Filled blocks — equivalent to moloch's `block=fill` theme option, but set as
% plain colors instead. The `block=fill` option is resolved through moloch's
% pgfkeys `.search also` forwarding, which only works in the newer moloch
% (TeX Live 2024); CI's older moloch (TeX Live 2023) rejects it with
% "I do not know the key '/block'". These four lines are what the option does.
\setbeamercolor{block title}{bg=normal text.bg!80!fg}
\setbeamercolor{block body}{bg=block title.bg!50!normal text.bg}
\setbeamercolor{block title alerted}{bg=block title.bg}
\setbeamercolor{block title example}{bg=block title.bg}
\setbeamertemplate{page number in head/foot}[totalframenumber]
\usepackage{amsmath, amssymb}
\usepackage{graphicx}
\usepackage{tikz}
\usepackage{comment}
\usepackage{hyperref}
\usepackage{subcaption}

\graphicspath{{fig/}}

\newcommand{\sectionframe}[1]{%
  \begin{frame}[plain]
    \centering
    \vfill
    {\Huge #1}
    \vfill
  \end{frame}
}

\title[]{Learning Dynamics}
\author[Corlouer]{Guillaume Corlouer}
\institute{Stormglass}
\date{}

\begin{document}
\maketitle
\begin{frame}{Big questions}
\begin{itemize}
    \item How do models generalize?
    \item How does internal structure form during training?
    \item Emergent capabilities
    \item Control levers for the training process
    \item How do the loss, architecture, data, and optimizer determine which internal structures are learned throughout training?
\end{itemize}
\end{frame}
\begin{frame}{How: making mathematical models of learning}
    \begin{itemize}
        \item Tractable toy models: deep linear networks, shallow nonlinear networks, tensor networks, polynomial activations
        \item Take limits in depth, width, data, and training time
        \item Exploit symmetries: quotient them out and find conserved quantities
        \item Identify useful observables: singular values of weights, learning coefficient, correlation functions
        \item Work in specific regimes: lazy versus rich
        \item Validate with tight feedback loops: experiments, symbolic derivations, and mathematical proofs
    \end{itemize}
\end{frame}
\begin{frame}{Safety motivations}
    \begin{itemize}
        \item A central question in AI safety is whether the training process selects a harmful model, such as a scheming, sycophantic, or spiteful model
        \item Predict whether a model will be harmful in high-stakes, out-of-distribution environments when behavioural or average-case interpretability evaluations cannot be trusted
        \item Develop stronger interpretability methods by decomposing models throughout training
        \item We want a mathematical model that can outperform average-case methods at predicting harmful behaviour
        \item More ambitiously, we would like to find control levers in the training process that make us confident training will select a safe model
    \end{itemize}
\end{frame}
\begin{frame}{Emergent misalignment}
  Jan Betley et al., ``Emergent Misalignment: Narrow finetuning can produce broadly misaligned LLMs'' (Figure 1)
\includegraphics[width=\textwidth,height=\textheight,keepaspectratio]{emergent_misalignment.png}
\end{frame}
\sectionframe{Setup}
\begin{frame}{Deep neural network}
\small
\begin{columns}[T,onlytextwidth]
\begin{column}{0.48\textwidth}
Let $\theta\in \Omega$ be the parameters, let $x\in\mathbb{R}^{d_0}$, and let $W_\ell\in\mathbb{R}^{d_\ell \times d_{\ell-1}}$ with biases $b_\ell$ and activation $\sigma$. With $h_0=x$,
\[
\begin{aligned}
h_\ell(\theta,x)
&=
\sigma\!\left(W_\ell h_{\ell-1}(\theta,x)+b_\ell\right), \\
&\hspace{1.5em}\ell=1,\ldots,L-1.
\end{aligned}
\]
A deep neural network is a function $f_\theta$ such that
\[
f_\theta(x) = W_Lh_{L-1}(\theta,x)
\]
Other architectural components include residual streams, tokenization, embeddings, attention for sequence modeling, and layer normalization.
\end{column}
\begin{column}{0.48\textwidth}
\centering
\vspace{0.2cm}
\resizebox{\linewidth}{!}{\input{fig/deep_neural_network.tex}}
\end{column}
\end{columns}
\end{frame}
\begin{frame}{Deep linear network}
The function is linear in the data and multi-linear (nonlinear) in the parameters:
  \[
    \begin{aligned}
      f\colon \mathbb{R}^{d_0} \times
      \prod_{\ell=1}^{L}\mathbb{R}^{d_\ell \times d_{\ell-1}}
      &\longrightarrow \mathbb{R}^{d_L}, \\
      (x,W_1,\ldots,W_L)
      &\longmapsto W_L\cdots W_1x = Wx.
    \end{aligned}
  \]
$d_0$ is the input dimension; $d_\ell$, $1\leq \ell \leq L$, is the width of layer $\ell$.
Define the product map:
\[
\mu(\theta) := W_L\cdots W_1 = W,
\]
where $W$ is the student matrix.
\end{frame}
\begin{frame}{Stochastic gradient descent}
\small
Learn a function $g$ from observations $y = g(x)$. Initialize neural weights at $\theta_0 \sim \mathcal{N}(0,\sigma I)$, with initialization scale $\gamma$ and $\sigma=d_L^{-\gamma}$.
\begin{block}{SGD update}
    Consider a dataset of size $N$. At each time step, sample a batch $b_k$ of size $B$. The empirical loss is the expectation of the batch loss:
\[
\mathcal{L}(\theta) = \mathbb{E}_b[\mathcal{L}(\theta; b)].
\]
    The SGD update is
\[
\begin{aligned}
\theta_{k+1}
&=
\theta_k-\eta\nabla_\theta\mathcal{L}(\theta_k;b_k) \\
&=
\theta_k-\eta\nabla_\theta\mathcal{L}(\theta_k)
+\eta\,\xi(\theta_k,b_k).
\end{aligned}
\]
with gradient noise
\[
\xi(\theta,b):=
\nabla_\theta \mathcal{L}(\theta)
-\nabla_\theta \mathcal{L}(\theta;b).
\]
\end{block}
The loss $\mathcal{L}$ can be, for example, mean squared error or cross-entropy loss.
\end{frame}
\begin{frame}{Data, loss function, gradient flow}
\begin{block}{Quadratic population loss and gradient flow}
Let $M$ be the teacher matrix, with $y=Mx$ and $x\sim \mathcal{N}(0,I_{d_0})$.
The input-output covariance is $\Sigma_{YX}=\mathbb{E}[yx^{\top}]=M$.
\begin{equation*}
    \mathcal{L}(\theta) := \frac{1}{2} \lVert M - W_L\cdots W_1\rVert^2_F = \frac{1}{2} \lVert M - W\rVert^2_F
\end{equation*}
Initialize $\theta_0 \sim \mathcal{N}(0,d_L^{-\gamma}I)$, where $\gamma$ controls the initialization scale. The gradient flow is
\[
\dot{W}_\ell = - \nabla_{W_{\ell}} \mathcal{L}(\theta).
\]
\end{block}
Frobenius norm: $\lVert A\rVert_F:=\sqrt{\operatorname{Tr}(A^{\top}A)}$.
\\
Global minimizers achieve zero loss: $\exists \theta^{\star}\in \Omega,\ \mu(\theta^{\star}) = W^{\star} = M$.
\\
The structure of the data is encoded by $\Sigma_{YX}$.
\end{frame}
\sectionframe{Geometry}
\begin{frame}[plain]
\centering
\includegraphics[width=\textwidth,height=\textheight,keepaspectratio]{global_minima_saddle_3d_loss.png}
\end{frame}
\begin{frame}{Deep learning without poor local minima}
Assume that $\Sigma_X$ and $\Sigma_{YX}$ are full rank, and that $\Sigma_{YX}$ has $d_L$ distinct singular values.
\begin{block}{No non-global local minima}
    The loss function of a DLN of depth $L$ is non-convex and non-concave, and every critical point that is not a global minimum is a saddle point.
\end{block}
The training dynamics cannot get stuck in spurious local minima.
\\
Reference: \textit{Kawaguchi, Kenji. ``Deep learning without poor local minima.'' Advances in neural information processing systems 29 (2016)}
\end{frame}
\begin{frame}{Parametrization of critical points}
Teacher SVD: \(M=\sum_{i=1}^{r_m}s_i u_i v_i^{\top}\).
\begin{block}{Saddle points}
    Let $\theta\in\Omega$ be a critical point of rank $r$. Then there exists a subset $S$ of left singular vectors of $M$ and a projector $P_S$ onto their span such that
    \[
    P_S := \sum_{i\in S}u_i u_i^{\top},
    \qquad
    \mu(\theta)=P_S M=\sum_{i\in S}s_i u_i v_i^{\top}.
    \]
\end{block}
The saddle points are given by a subset of singular values and (left) singular vectors of the teacher $M$.
\\
Reference: \textit{Achour, El Mehdi, François Malgouyres, and Sébastien Gerchinovitz. ``The loss landscape of deep linear neural networks: a second-order analysis.'' Journal of Machine Learning Research 25.242 (2024): 1-76.}
\end{frame}
\begin{frame}{Symmetries}
\small
Let $G_{\mathbf{d}}:=\prod_{\ell=0}^{L} \operatorname{GL}_{d_\ell}$ act on $\Omega$ via
\[
(W_1,\ldots,W_L)
\mapsto
(g_1W_1g_0^{-1},\; g_2W_2g_1^{-1},\;\ldots,\; g_L W_L g_{L-1}^{-1}).
\]
\begin{block}{Equivariance of the multiplication map}
    Since \(W_Lg_{L-1}^{-1}g_{L-1}W_{L-1}\cdots g_1^{-1}g_1W_1=W_L\cdots W_1\), the multiplication map $\mu$ is $G_{\mathbf{d}}$-equivariant and invariant under the hidden gauge group $G_h$:
    \[
    \mu(g\cdot \theta)=g_L\mu(\theta)g_0^{-1},
    \qquad
    \mu(g_h\cdot \theta) = \mu(\theta).
    \]
A DLN is degenerate: infinitely many parameter values can map to the same student matrix.
\end{block}
It is natural to study solutions up to $G_{\mathbf{d}}$ symmetries. By equivariance of $\mu$, $G_{\mathbf{d}}$ acts on the determinantal variety
\[
\Sigma_{\mathbf{d}}^{r}:=\{\theta\in \Omega\mid\operatorname{rank}(\mu(\theta)) = r\}.
\]
\end{frame}
\begin{frame}{Stratification of the fiber of the product map}
\begin{block}{Orbits are parameterized by rank patterns}
The $G_{\mathbf{d}}$-orbits in $\Sigma_{\mathbf{d}}^r$ are indexed by rank patterns
\[
\rho_{ij}:=\operatorname{rank}(W_j\cdots W_i),
\qquad
1\leq i\leq j\leq L.
\]
The corresponding orbit decomposition is
\[
\Sigma_{\mathbf{d}}^r = \bigsqcup_{\rho} \mathcal{O}_{\rho}.
\]
Two DLNs are in the same $G_{\mathbf{d}}$-orbit if and only if they have the same rank pattern.
\\
Since $\Sigma_{\mathbf{d}}^{r} = G_{\mathrm{out}}\cdot\mu^{-1}(W)$, this gives a combinatorial description of the fiber of the product map $\mu^{-1}(W)$.
\end{block}
Rank patterns are natural observables arising from geometric invariants of DLNs.
\end{frame}
\begin{frame}{Link with SLT}
\begin{block}{Codimension of the global minima}
For a teacher matrix $M$ of rank $r_m$, we have:
    \[
    \begin{aligned}
    \operatorname{codim}\bigl(\mu^{-1}(M)\bigr)
    &=
    \operatorname{codim}\bigl(\Sigma_{\mathbf{d}}^{r_m}\bigr)
    + r_m(d_0 + d_L - r_m) \\
    &= 2\,\operatorname{RLCT}_{0}(\mathcal{L}).
    \end{aligned}
    \]
In general, we only have:
\[
\operatorname{RLCT}_{0}(\mathcal{L})
\leq
\frac{1}{2}\operatorname{codim}_{\operatorname{Rep}_{\mathbf{d}}}\bigl(\mathcal{L}^{-1}(0)\bigr).
\]
\end{block}
Ref: \textit{Lehalleur, Simon Pepin, and Richárd Rimányi. ``Geometry of fibers of the multiplication map of deep linear neural networks.'' arXiv preprint arXiv:2411.19920 (2024).}
\end{frame}
\begin{frame}{Discussion}
\begin{block}{Takeaways}
Under full-rank and nondegenerate-data assumptions:
    \begin{itemize}
        \item Critical points of DLNs are either saddles or global minima
        \item Gauge symmetries mean there is a whole algebraic variety of degenerate minima
        \item Geometric invariants give observables and a combinatorial description of the parameter-function map and the global minima
    \end{itemize}
\end{block}
\end{frame}
\begin{frame}{References}
\begin{itemize}
    \item Kenji Kawaguchi. Deep learning without poor local minima. \textit{Advances in Neural
Information Processing Systems}, 29, 2016.
    \item El Mehdi Achour, François Malgouyres, and Sébastien Gerchinovitz. The loss
landscape of deep linear neural networks: a second-order analysis. \textit{Journal of
Machine Learning Research}, 25(242):1–76, 2024.
    \item Anna Choromanska, Mikael Henaff, Michael Mathieu, Gérard Ben Arous, and
Yann LeCun. The loss surfaces of multilayer networks. In \textit{Artificial Intelligence
and Statistics}, pages 192–204. PMLR, 2015.
    \item Lehalleur, Simon Pepin and Rim{\'a}nyi, Rich{\'a}rd. Geometry of fibers of the multiplication map of deep linear neural networks, arXiv preprint arXiv:2411.19920, 2024.
    \item Chen, Po and Jiang, Rujun and Wang, Peng. A Complete Loss Landscape Analysis of Regularized Deep Matrix Factorization, arXiv preprint arXiv:2506.20344, 2025.
\end{itemize}
\end{frame}
\sectionframe{Dynamics}
\begin{frame}{Gradient flow and NTK dynamics}
  \small
  In the small-learning-rate limit, gradient descent becomes gradient flow:
  \[
    \theta_{t+1}=\theta_t-\eta \nabla_\theta \mathcal{L}(\theta_t)
    \qquad \leadsto \qquad
    \dot{\theta}(t)=-\nabla_\theta \mathcal{L}(\theta(t))
  \]

  Neural Tangent Kernel (NTK):
  \[
    K_{\theta}(x,x')
    =
    \nabla_{\theta} f_{\theta}(x)\, \cdot
    \nabla_{\theta} f_{\theta}(x')^{\top}
    =
    \sum_{i=1}^p
    \partial_{\theta_i} f_{\theta}(x)\,
    \partial_{\theta_i} f_{\theta}(x')^{\top}.
  \]

  For a quadratic loss on a dataset $(x_i,y_i)$, the dynamics in function space are
  \[
    \frac{d}{dt}f_{\theta_t}(x)
    =
    -\sum_{i=1}^N
    K_{\theta_t}(x,x_i)\bigl(f_{\theta_t}(x_i)-y_i\bigr).
  \]
\end{frame}
\begin{frame}{Conserved quantities in DLNs}
  Conserved quantities under gradient flow:
  \[
    \beta_\ell :=
    W_{\ell+1}^{\top}W_{\ell+1}
    -
    W_{\ell}W_{\ell}^{\top},
    \qquad
    \ell=1,\ldots,L-1.
  \]

  Along the flow, $\beta_\ell$ is constant. When $\beta_\ell=0$, we have \textbf{balanced weights}.

  \vspace{0.4cm}

  For a one-dimensional, two-layer DLN, we get a hyperbolic equation:
  \[
  w_2^2-w_1^2=\beta.
  \]

  \vspace{0.4cm}
\end{frame}
\begin{frame}[plain]
\centering
\includegraphics[width=\textwidth,height=\textheight,keepaspectratio]{foliation_paramaters_balanced.png}
\end{frame}
\begin{frame}{The manifold of gradient flow}
\small
Gradient flow is constrained to the following hyperbolic manifold:
\begin{equation*}
    \mathcal{M}_{\beta}
    :=
    \{\theta\in\Omega\mid
    W_{\ell+1}^{\top}W_{\ell+1}
    -
    W_{\ell}W_{\ell}^{\top}
    =
    \beta_\ell,\;
    \ell=1,\ldots,L-1\}.
\end{equation*}
\begin{block}{Gradient flow on the balanced variety}
On the balanced variety $\mathcal{M}_0$, the self-consistent gradient flow equation in function space is
\begin{equation*}
    \dot{W}
    =
    \sum_{k=1}^{L}
    (WW^{\top})^{\frac{L-k}{L}}
    R(W)
    (W^{\top}W)^{\frac{k-1}{L}},
    \qquad
    R(W):=\Sigma_{YX}-W.
\end{equation*}
For depth $L=2$:
\[
    \mathcal{M}_0
    =
    \{(W_1,W_2)\mid W_2^{\top}W_2=W_1W_1^{\top}\}.
\]
\begin{equation*}
    \dot{W}
    =
    (WW^{\top})^{\frac{1}{2}}R(W)
    +
    R(W)(W^{\top}W)^{\frac{1}{2}}.
\end{equation*}
\end{block}
\end{frame}
\begin{frame}{Learning rate: Edge of stability}
Descent lemma: for an $L_{\mathrm{smooth}}$-smooth convex loss, the loss decreases if $0<\eta<2/L_{\mathrm{smooth}}$.
\\
The loss is nonconvex, but the Hessian oscillates near and slightly above $\frac{2}{\eta}$; this regularizes away sharp minima.

  \centering
  \includegraphics[width=\textwidth,height=0.57\textheight,keepaspectratio]{edge_of_stability.png}


  \begin{minipage}{0.95\textwidth}
    \centering
    Ref: \href{https://arxiv.org/abs/2103.00065}{Gradient Descent on Neural Networks Typically Occurs at the Edge of Stability, Cohen et al.}
  \end{minipage}
\end{frame}
\begin{frame}{Takeaways}
  \begin{block}{Key points}
    \begin{itemize}
      \item NTK preconditions gradient flow in function space
      \item Symmetries can induce conserved quantities that constrain the gradient flow
      \item Gradient flow in function space for balanced quantities leads to self-consistent equations
      \item The effect of the learning rate can be understood via the edge of stability
    \end{itemize}
  \end{block}
\end{frame}
\begin{frame}{References}
    \begin{itemize}
        \item Menon, Govind. The geometry of the deep linear network. Symposium on Probability and Stochastic Processes. Cham: Springer Nature Switzerland, 2023.
        \item Sibylle Marcotte, R\'emi Gribonval, and Gabriel Peyr\'e. Abide by the law and
follow the flow: Conservation laws for gradient flows. \textit{Advances in neural
information processing systems}, 36:63210–63221, 2023.
    \item Sanjeev Arora, Nadav Cohen, Noah Golowich, and Wei Hu. \textit{A convergence
analysis of gradient descent for deep linear neural networks}. arXiv preprint
arXiv:1810.02281, 2018
    \item Arthur Jacot, Franck Gabriel, and Cl\'ement Hongler. Neural tangent kernel:
    Convergence and generalization in neural networks. \textit{Advances in Neural Information
    Processing Systems}, 31, 2018.
    \item Jeremy M. Cohen, Alex Damian, Ameet Talwalkar, J. Zico Kolter, and Jason D.
    Lee. Understanding optimization in deep learning with central flows. arXiv
    preprint arXiv:2410.24206, 2024.
    \end{itemize}
\end{frame}
\sectionframe{Initialization, Width, and Depth}
\begin{frame}{Training losses}
Initialize weights: $\theta \sim \mathcal{N}(0,d_L^{-\gamma}I)$
\\
$\gamma < 1$: large initialization, linear regime (NTK)
\\
$\gamma > 1$: small initialization, saddle-to-saddle
\\
$\gamma = 1$: intermediate regime

    \begin{figure}
        \centering
        \includegraphics[width=\linewidth,height=0.56\textheight,keepaspectratio]{DLN_losses_regimes.png}
    \end{figure}
\end{frame}

\begin{frame}{Singular values of the student matrix}
    \begin{figure}
        \centering
        \includegraphics[width=\linewidth,height=0.75\textheight,keepaspectratio]{DLN_sv_students.png}
    \end{figure}
\end{frame}

\begin{frame}{Singular values of individual weight matrices}
    \begin{figure}
        \centering
        \includegraphics[width=\linewidth,height=0.75\textheight,keepaspectratio]{DLN_individual_sv.png}
    \end{figure}
\end{frame}
\begin{frame}{Lazy regime (NTK or linear regime)}
  \small
  For a DLN of depth $L$ with initialization $\theta(0)\sim \mathcal{N}(0,\sigma^2 I)$ and $\sigma^2 = d_L^{-\gamma}$, with $\gamma < 1$:
  \[
    \dot{W} \approx K_0(M-W).
  \]

  This leads to the following solution:
  \[
    W(t)=M-\exp(-K_0t)\bigl(M-W(0)\bigr).
  \]

  The timescale of learning is fixed by the NTK at initialization:
  \[
    t \approx -K_0^{-1}\log\left(\frac{s - w_f}{s - w_0}\right)
  \]

  This behaves like a shallow linear network.
\end{frame}
\begin{frame}{Rich regime}
  \small
  For balanced, and aligned weights, and $\gamma > 1$ (small init.), the singular values of the student satisfy
  \[
    \dot{w}_{\alpha} = 2\left(s_{\alpha} - w_{\alpha}\right)w_{\alpha}.
  \]

  Solving this ODE for $w(t):=w_{\alpha}$ gives:
  \[
    w(t)=\frac{s}{1+\left(\frac{s}{w_0}-1\right)e^{-2s t}}.
  \]

  The timescale of learning is
  \[
    t=\frac{1}{2s}\,\ln\!\left(\frac{w_f\,(s-w_0)}{w_0\,(s-w_f)}\right), \qquad s\neq 0.
  \]
\end{frame}
\begin{frame}{Takeaways}
  \begin{block}{Key points}
    \begin{itemize}
      \item Depth and specific initialization induce saddle-to-saddle dynamics
      \item Gradient flow has a simplicity bias toward learning low hanging fruit directions
      \item Large initialization and large width give the NTK regime: no feature learning, and the NTK is frozen
      \item In the saddle-to-saddle regime, the timescale of learning is set by the target singular values
    \end{itemize}
  \end{block}
\end{frame}
\begin{frame}{References}
    \begin{itemize}
        \item Saxe, Andrew M., James L. McClelland, and Surya Ganguli. A mathematical theory of semantic development in deep neural networks. \textit{Proceedings of the National Academy of Sciences} 116.23 (2019): 11537-11546.
        \item Jacot, Arthur, et al. Saddle-to-saddle dynamics in deep linear networks: Small initialization training, symmetry, and sparsity. \textit{arXiv preprint} arXiv:2106.15933 (2021).
        \item Dominé, Clémentine CJ, et al. From lazy to rich: Exact learning dynamics in deep linear networks. \textit{arXiv preprint} arXiv:2409.14623 (2024).
        \item Zhang, Yedi, Andrew Saxe, and Peter E. Latham. Saddle-to-Saddle Dynamics Explains A Simplicity Bias Across Neural Network Architectures. \textit{arXiv preprint} arXiv:2512.20607 (2025).
        \item Zhenfeng Tu, Santiago Tomas Aranguri Diaz, and Arthur Jacot. Mixed dynamics
in linear networks: Unifying the lazy and active regimes. \textit{Advances in Neural
Information Processing Systems}, 37:106059–106104, 2024.
    \end{itemize}
\end{frame}
\sectionframe{Stochasticity}

\begin{frame}{Langevin dynamics}
  \begin{itemize}
    \item SGD: Sample a batch from a dataset and compute the batch gradient
    \item Continuous model with Brownian motion $B_t$
  \end{itemize}

  \[
    d\theta_t
    =
    -\nabla \mathcal{L}(\theta_t)\,dt
    +
    \sqrt{\eta}\,\Sigma^{1/2}(\theta_t)\,dB_t.
  \]

  \begin{itemize}
    \item The gradient-noise covariance $\Sigma(\theta_t)$ depends on geometry: it is anisotropic and state-dependent
  \end{itemize}
\end{frame}

\begin{frame}{Fokker-Planck: Boltzmann equilibria}
  \small
  The Fokker-Planck equation (for density $p(\theta,t)$) is a local conservation law:
  \[
    \begin{aligned}
      \partial_t p(\theta,t) &= -\nabla\!\cdot j(\theta,t), \\
      j(\theta,t)
      &=
      -p(\theta,t)\nabla \mathcal{L}(\theta)
      -\frac{\eta}{2}\nabla_\theta^{\top}\!\left(\Sigma(\theta)p(\theta,t)\right).
    \end{aligned}
  \]

  Important special case:
  \begin{itemize}
    \item Stationary distribution: $\partial_t p^{\ast}(\theta)=0$
    \item Thermal equilibrium: $j=0$
    \item White noise: $\Sigma=\sigma^2 I$
  \end{itemize}

  \[
    p^{\ast}(\theta)
    \propto
    \exp\!\left(-\frac{2}{\eta\sigma^2}\mathcal{L}_N(\theta)\right)
    \qquad \text{(Boltzmann)}.
  \]
\end{frame}

\begin{frame}{Noise anisotropy}
  \small
  \begin{itemize}
    \item In general, noise is anisotropic: it introduces a different implicit bias
    \item In function space, Langevin dynamics introduce an additional drift
  \end{itemize}

  \[
    \frac{\eta}{2}\operatorname{Tr}\!\bigl(\Sigma(\theta)H(\theta)\bigr)
  \]

  \begin{itemize}
    \item Heuristic: bias toward flatter minima
    \item More:
      \begin{itemize}
        \item See Govind Menon \href{https://arxiv.org/abs/2602.12257}{On the implicit regularization of Langevin dynamics with projected noise}
        \item Our \href{https://arxiv.org/html/2604.06366v1}{paper} deriving Langevin dynamics for DLNs during the saddle-to-saddle regime
      \end{itemize}
  \end{itemize}
\end{frame}

\begin{frame}{Golden Path hypothesis}
  \small
  \begin{itemize}
    \item Lots of debate about the implicit bias of SGD noise
    \item Transient regime (drift-dominated) versus low-loss regime (diffusion-dominated):
  \end{itemize}

  \centering
  \includegraphics[height=0.6\textheight]{golden_path.png}

  \vspace{0.15cm}
  \footnotesize
  \emph{Beyond Implicit Bias: The Insignificance of SGD Noise in Online Learning}
\end{frame}
\sectionframe{Towards nonlinearity}

\begin{frame}{Neural race reduction in gated DLNs}
  \centering
  \includegraphics[height=0.68\textheight]{neural_race.png}

  \vspace{0.2cm}
  \footnotesize
  \emph{The Neural Race Reduction: Dynamics of Abstraction in Gated Networks, Andrew M. Saxe et al.}
\end{frame}

\begin{frame}{MLP and linear transformers}
  \small
  \emph{Saddle-to-Saddle Dynamics Explains A Simplicity Bias Across Neural Network Architectures, Saxe et al.}

  \vspace{0.2cm}
  \centering
  \includegraphics[height=0.58\textheight]{simplicity_bias_transformers_mlp.png}
\end{frame}
\begin{frame}{Alternating Gradient Flows}
  \small
  \emph{Alternating Gradient Flows: A Theory of
Feature Learning in Two-layer Neural Networks}
  \vspace{0.2cm}
  \centering
  \includegraphics[height=0.58\textheight]{AF_architectures.png}
\end{frame}
\begin{frame}
  \small
  \emph{Alternating Gradient Flows: A Theory of
Feature Learning in Two-layer Neural Networks}
  \vspace{0.2cm}
  \centering
  \includegraphics[height=0.58\textheight]{AF_conceptual.png}
\end{frame}
\begin{frame}{Takeaways}
  \begin{block}{Key points}
    \begin{itemize}
      \item SGD can be modeled by an SDE with state-dependent noise
      \item If the noise is white and the system is at thermal equilibrium, the stationary distribution of the SDE is Boltzmann
      \item The implicit bias of stochasticity is not fully understood (flatness?), but it is more of a second-order effect
      \item Nonlinearity can be modeled with gated DLNs: neural race reduction
      \item MLPs and linear transformers also have saddle-to-saddle dynamics
      \item Alternating gradient flow: towards a general theory of the rich regime
    \end{itemize}
  \end{block}
\end{frame}
\begin{frame}{Promising directions}
\begin{itemize}
    \item Polynomial neural networks (quadratic activations)
    \item Related: tensor networks and tree tensor networks
\end{itemize}
\begin{block}{Big questions}
\begin{itemize}
  \item Toy models of deep nonlinear learning
  \item Theory that captures natural data
  \item Making explicit the implicit biases of realistic architectures
  \item Mixing rich and lazy regimes
  \item A statistical field theory of deep neural networks
\end{itemize}
\end{block}
Reference: \href{https://arxiv.org/abs/2604.21691}{\textit{There Will Be a Scientific Theory of Deep Learning, Jamie Simon et al.} }
\end{frame}
\end{document}
